\section{本轮收敛诊断与可复核终态}
2026 年 10 月的补充运行采用统一任务接口，分别保存执行状态、物理验收、源码与输入身份。制造解在 208、756 和 1856 个节点上通过；最大场值误差为 $1.70\times10^{-14}$，最大导数误差为 $2.26\times10^{-13}$。
圆柱的解析边界迹与二维离散解边界迹在相同的 1683、3672 和 9108 节点挤出网格上配对重算。最细档解析迹的 $J$ 相对误差为 0.877248\%，$\phi$ 误差为 $1.70939\times10^{-4}$，gauge 残差为 $2.77686\times10^{-2}$；离散迹对应值为 0.877132\%、$1.47596\times10^{-4}$ 和 $2.77850\times10^{-2}$。两组均保留实际 $\phi$ 及其导数。仅替换边界迹未使 $\phi$ 和 gauge 达到 $10^{-8}$ 与 $10^{-3}$ 门槛；下一实验应检查挤出离散算子的分块残差与 gauge 耦合。

Maxwell 七股采用官方网格链接并完成一次有依据的设置修订。源设计和基解均为 1694 个实际网格单元，但严格坐标身份仍失败。仅允许节点和单元编号重排的规范化，未放宽坐标门槛。因此，本轮交付固定网格诊断；后续网格收敛、边界敏感性与同条件跨求解器比较保持依赖阻塞状态。历史单网格一致性继续按其限定范围引用。

机器可读证据索引位于仓库的 provenance/convergence 目录。上述有限终态支持后续最小实验设计；最终 Q2 设计仍要求名义铜面积至少 $5\,\mathrm{mm^2}$、传统多级 Litz 强基线以及实际弧长和铜量公平比较。
